\MWChapter{Method, provenance, and reproducibility}{方法、溯源与可复现性}
\label{ch:methods}

\section{Evidence model / 证据模型}

The relevant reproducibility unit is not a dataset split or model checkpoint. It is a chain that
binds a named theorem to an immutable source revision, a build/audit record, and a manifest entry.
The current delivery record identifies source \(S\), evidence \(E\), and a later pointer/manifest
revision \(P\); cleanup or documentation commits do not replace the proof identity of that chain
\citep{cantiluneClosure}.

\begin{table}[H]
\centering
\caption{Evidence fields required to interpret a formal claim.}
\label{tab:manifest}
\begin{tabularx}{\textwidth}{@{}L{0.26\textwidth}Y L{0.26\textwidth}@{}}
\toprule
\rowcolor{MWSoft!70}
\textbf{Field} & \textbf{Current record} & \textbf{Why it matters} \\
\midrule
Source branch & \texttt{codex/theory-foundation} & Fixes the reviewed repository state \\
Verified source \(S\) & \texttt{59a1a688\ldots e67d15dc} & Identifies the Lean source tree behind the stated proof status \\
Evidence \(E\) & \texttt{ed26cb74\ldots 3fb8dee5} & Binds the build/audit record to the source claim \\
Manifest & \texttt{formal/proof-obligations.json} & Names declarations, status, provenance, and review fields \\
Build evidence & QA record bound to source \(S\) & Records the source-integrity and audit outcome \\
Review status & \emph{review-pending} & Prevents a local or automated build from being reported as independent approval \\
\bottomrule
\end{tabularx}
\end{table}

\section{Rebuild protocol / 重建协议}

The formal project pins Lean and mathlib through its toolchain and lock files. Routine work must
not silently update dependencies; the repository documents an ordinary kernel rebuild and a
separate provenance-aware evidence gate \citep{cantiluneFormalReadme}.

\begin{MWCodeBox}[Repository verification commands]
\begin{lstlisting}[style=moonweave,language=bash]
lake build
.\formal\scripts\ci.ps1
.\formal\scripts\ci.ps1 -RequireProved -VerifyTreeOnly
\end{lstlisting}
\end{MWCodeBox}

The commands are not an installation guide for an end-user product. The first rebuilds the
kernel-checked project with the pinned toolchain; the evidence scripts verify source integrity,
known dependencies, and the relationship between declared proof status and durable evidence.

\section{Methodological exclusions / 方法上的排除项}

No training corpus, benchmark dataset, model provider revision, inference hardware profile, or
latency/cost experiment is part of the present evidence package. Such artifacts would become
necessary only for a future product or runtime report. Reporting no empirical result is more
accurate than repurposing template curves or synthetic measurements.

\begin{MWRisk}
The status \emph{proved / review-pending} is intentionally narrower than “validated product.”
A successful local build is neither an independent review nor a substitute for FCP disposition,
ADR acceptance, or per-product conformance evidence.
\end{MWRisk}
